
\begin{UPSTIexercice}{Un \texttt{switch} bien construit}
    On exécute le programme suivant, avec \texttt{niveau = 2}.
    \begin{lstlisting}[language=c]
int niveau = 2;

switch (niveau)
{
    case 1:
        printf("Debutant\n");
        break;
    case 2:
        printf("Intermediaire\n");
        break;
    case 3:
        printf("Expert\n");
        break;
    default:
        printf("Niveau inconnu\n");
        break;
}
printf("Fin du programme\n");
    \end{lstlisting}
    \UPSTIquestion{Prédire l'affichage de ce programme.}
    \UPSTIquestion{Le \texttt{case 3} et le \texttt{default} sont-ils exécutés ?}
\end{UPSTIexercice}

\begin{UPSTIprofOnlyEnv}
    \begin{UPSTIcorrectionP}{Un \texttt{switch} bien construit}
        Affichage attendu~:
        \begin{lstlisting}[language=bash,style=console]
Intermediaire
Fin du programme
        \end{lstlisting}
        \texttt{niveau} vaut \texttt{2} : seul le \texttt{case 2} correspond, il affiche \og Intermediaire \fg{}, puis le \texttt{break} interrompt immédiatement le \texttt{switch}. Ni \texttt{case 3} ni \texttt{default} ne sont exécutés, car un \texttt{switch} ne teste qu'un seul \texttt{case} correspondant (grâce au \texttt{break}) avant de poursuivre après l'accolade fermante.
    \end{UPSTIcorrectionP}
\end{UPSTIprofOnlyEnv}


\begin{UPSTIexercice}{Piège classique : le \texttt{break} oublié}
    On exécute le programme suivant, avec \texttt{jour = 3}.
    \begin{lstlisting}[language=c]
int jour = 3;

switch (jour)
{
    case 1:
        printf("Lundi\n");
    case 2:
        printf("Mardi\n");
    case 3:
        printf("Mercredi\n");
    case 4:
        printf("Jeudi\n");
        break;
    case 5:
        printf("Vendredi\n");
        break;
}
    \end{lstlisting}
    \UPSTIquestion{Prédire ce qui s'affiche dans la console. \textbf{Attention}, ce programme comporte une erreur classique.}
    \UPSTIquestion{Nommer ce comportement, et expliquer en une phrase pourquoi il se produit.}
    \UPSTIquestion{Comment corriger le programme pour qu'il affiche uniquement \og Mercredi \fg{} ?}
\end{UPSTIexercice}

\begin{UPSTIprofOnlyEnv}
    \begin{UPSTIcorrectionP}{Piège n°1 : le \texttt{break} oublié}
        Affichage attendu~:
        \begin{lstlisting}[language=bash,style=console]
Mercredi
Jeudi
        \end{lstlisting}
        \texttt{jour} vaut \texttt{3} : le \texttt{case 3} correspond et affiche \og Mercredi \fg. Comme il n'y a \textbf{pas de \texttt{break}} à la fin de ce \texttt{case}, l'exécution \textbf{continue} dans le \texttt{case} suivant (\texttt{case 4}), qui affiche \og Jeudi \fg, jusqu'à rencontrer le premier \texttt{break}. Ce comportement s'appelle le \textbf{fall-through} : en l'absence de \texttt{break}, un \texttt{switch} ne s'arrête pas à la fin d'un \texttt{case}, il \og tombe \fg dans les suivants.

        Pour n'afficher que \og Mercredi \fg, il suffit d'ajouter un \texttt{break;} juste après la ligne \texttt{printf("Mercredi\textbackslash n");} du \texttt{case 3}.
    \end{UPSTIcorrectionP}
\end{UPSTIprofOnlyEnv}

\pagebreak
\begin{UPSTIexercice}{Fall-through volontaire}
    On exécute le programme suivant. On rappelle qu'un jour est \og ouvré \fg s'il n'est ni samedi ni dimanche. La variable \texttt{jour} vaut \texttt{6} (1 = lundi, ..., 7 = dimanche).
    \begin{lstlisting}[language=c]
int jour = 6;

switch (jour)
{
    case 6:
    case 7:
        printf("Week-end\n");
        break;
    default:
        printf("Jour ouvre\n");
        break;
}
    \end{lstlisting}
    \UPSTIquestion{Prédire l'affichage de ce programme.}
\end{UPSTIexercice}

\begin{UPSTIprofOnlyEnv}
    \begin{UPSTIcorrectionP}{Piège n°2 : fall-through volontaire et \texttt{default} mal placé}
        Affichage obtenu~: \texttt{Week-end}.

        Le \texttt{switch} compare \texttt{jour} à \textbf{chaque} \texttt{case}, dans l'ordre du code, et non le \texttt{default} en premier systématiquement : \texttt{default} n'est retenu que si \textbf{aucun} \texttt{case} ne correspond à la valeur testée. Ici, \texttt{jour = 6} correspond au \texttt{case 6}, donc c'est ce bloc (partagé avec \texttt{case 7} par fall-through) qui s'exécute, quelle que soit la position du \texttt{default} dans le fichier.

        Il s'agit bien ici d'un fall-through \textbf{volontaire} : les deux \texttt{case} (\texttt{6} et \texttt{7}) sont \textbf{vides}, c'est-à-dire qu'aucune instruction ne les sépare. Cette écriture est une façon standard de dire \og pour \texttt{jour} valant \texttt{6} \textbf{ou} \texttt{7}, exécuter le même bloc \fg, sans dupliquer le code. On la distingue d'un oubli (exercice précédent) au fait qu'elle est cohérente avec l'intention du programme et ne casse pas la logique attendue.
    \end{UPSTIcorrectionP}
\end{UPSTIprofOnlyEnv}
